\documentclass{article}
\usepackage{../math_template}

\author{Jonathan Kasongo}
\title{Polish Math Olympiad 2007 Round 1 --- P6/12}

\begin{document}
\maketitle

\begin{problem}{Polish Math Olympiad 2007 Round 1 --- P6/12}
Find all polynomials $W(x)$ such that $W\left(x^2\right) W\left(x^3\right)
= W(x)^5$ holds for every real number $x$.
\end{problem}

\begin{solution}{Solution}
Let the roots of $W$ be $|r_1| \leq |r_2| \leq \cdots \leq |r_n|$ where
$(r_n) \in \mathbb{C}$. The equation looks like so,

$$
\left(x^2 - r_1\right)\left(x^3 - r_1\right)\cdots
\left(x^2 - r_n\right)\left(x^3 - r_n\right)
=
(x - r_1)^5\cdots(x - r_n)^5
$$

and so the set of roots of both sides of this equation must be the same.\\

\textit{Claim 1:} If $r_i$ is a root of $W$, then $\pm \sqrt{r_i}$ and
$\sqrt[3]{r_i}$ are roots of $W$ aswell.\\

\textit{Proof:} On the LHS of the above equation we will find the factors
$\left(x^2 - r_i\right)\left(x^3 - r_i\right)$ which have solutions
$\pm \sqrt{r_i}$ and $\sqrt[3]{r_i}$. $\square$\\

\textit{Claim 2:} $r_1 = r_2 = \cdots = r_n = 0$. \\

\textit{Proof:}
Suppose, for the sake of contradiction, that there exists
an integer $k$ such that $r_k \neq 0$. \\

In fact let $k^{*}$ be the
smallest such integer for which $r_{k^{*}} \neq 0$. Now since
$\sqrt{r_{k^{*}}} \neq 0$ is a root of the LHS, it is also a root of the
RHS. But since $\sqrt{r_{k^{*}}} \leq r_{k^{*}}$ and $k^{*}$ is the smallest
integer for which $r_{k^{*}} \neq 0$, it must be the case that
$\sqrt{r_{k^{*}}} = r_{k^{*}}$. Since it cannot be 0, we have
$r_{k^{*}} = 1$. Now consider the following, all due to claim 1:

\[
\begin{aligned}
1 \text{ is a root of } W &\implies
-\sqrt{1} = -1 \text{ is a root of } W &\implies \\
\sqrt{-1} = i \text{ is a root of } W &\implies
\sqrt{i} = e^{i\pi/2} \text{ is a root of } W &\implies \\
\sqrt{e^{i\pi/2}} = e^{i\pi/4} \text{ is a root of } W &\implies
\sqrt{e^{i\pi/4}} = e^{i\pi/8} \text{ is a root of } W &\implies \\
\cdots &\implies \cdots &\implies
\end{aligned}
\]

and thus $e^{i\pi/2^n}$ is a root of $W$ for any positive integer $n$.
However since a polynomial can't have infinite roots, we have reached a
contradiction. $\square$ \\

So since $r_1 = r_2 = \cdots = r_n = 0$, it must be the case that
$W(x) = x^n$ for positive integer $n$, and this works since
$x^{2n} \cdot x^{3n} = (x^n)^5$. $\blacksquare$
\end{solution}
\end{document}
